% Generated by tools/build_v6.py; edit the assembly rules, not this file.
\section{Rare-event localization and margins below two}\label{app:localization}
This appendix supplies the general lower bound, its optimizer, the
below-two matching construction, and a general upper bracket.
The earlier arbitrary-law conjecture has been replaced by its correction;
\cref{sec:entropy-tail} gives the arbitrary-law characterization.
The published rare-prior input is identified explicitly below.

\subsection{Experiment, localization, and the two anchors}
A fresh coin has $D\sim\mu$ on $(0,1/2]$, $S\mid D=d\sim\operatorname{Ber}(q(d))$, and probability
\[
 P=\tfrac12+(2S-1)D.
\]
Conditional on its latent probability, successive observations of that coin are independent Bernoulli trials. Distinct fresh coins are independent. An algorithm can create coins, revisit them, interleave their observations, and stop adaptively. It observes neither $D$ nor $S$. One observation is one unit of work. Arithmetic and independent randomization are not charged, as in the original model.

Fix $0<\rho\le1/4$ and define
\begin{equation}\label{loc:eq:class}
 \Q_\rho(\mu)=\left\{q:(0,1/2]\to[0,1]\text{ measurable}:\quad
 \theta(q):=\int q\,d\mu\le\rho\right\}.
\end{equation}
The quantity to estimate is $\theta(q)$, not the known upper bound $\rho$. Replacing the inequality by the known equality $\theta=\rho$ would trivialize estimation.

Let $R_{\mu,\rho}(B)$ be the infimum, over algorithms with $\sup_{q\in\Q_\rho}\E_qT\le B$, of $\sup_q(\E_q|\widehat\theta-\theta(q)|^2)^{1/2}$. Let $r_{\mu,\rho}(B,\delta)$ be the corresponding smallest uniform error radius at failure probability at most $\delta$. Here $B\ge1$ and $0<\delta\le1/4$. These are \emph{expected-work} criteria. The minimax benchmark may know $\mu,\rho$. Each construction's actual parameter requirements are stated separately.

This class deliberately permits the positives to concentrate in a difficult subpopulation. It does \emph{not} impose the stronger condition $q(d)\le\rho$ at every magnitude. Nor does it assume known conditional magnitude laws $\mu_+,\mu_-$. In the latter model the mixture weight is the main unknown; in \eqref{loc:eq:class}, changing $q$ changes both conditional laws. Lee--Valiant's Section 7 addresses the known-conditional-law experiment with a per-coin depth restriction~\cite{leevaliant}.

\subsubsection{Sanity checks}
For $\mu=\delta_\Delta$, the problem is the fixed-gap mixture experiment. The Lee--Valiant bounds~\cite{leevaliant}, localized to $\theta\le\rho$, give, in the nontrivial regime $\varepsilon\le c\rho$,
\begin{equation}\label{loc:eq:lvanchor}
 B^*(\varepsilon,\delta)\asymp
 \frac{\rho}{\varepsilon^2\Delta^2}\log(1/\delta).
\end{equation}
The corresponding error scale, saturated at the trivial radius, is
\[\min\{\rho,\sqrt{\rho\log(1/\delta)/B}/\Delta\}.\]
At fixed confidence or for RMS, omit the logarithmic factor. This is a cited anchor, not a new fixed-gap theorem.

For constant $\rho$, a correct localized scale must agree, up to constants, with the original
\[
 \omega_\mu(h)=\E_\mu\min\{1,h/D\},\qquad h=B^{-1/2}
\]
(and $h=\sqrt{\log(1/\delta)/B}$ for high confidence). The proposed heuristic
\begin{equation}\label{loc:eq:heuristic}
 H_{\mu,\rho}(h)=\rho\,\omega_\mu(h/\sqrt\rho)
\end{equation}
satisfies both anchors. The next section shows that this is insufficient.

\Needspace{8\baselineskip}
\subsection{A counterexample to the proposed rescaling}
\begin{theorem}\label{loc:thm:counter}
Let $0<\rho\le1/4$, $a=1/4$, $d=\rho^2$,
\[
 \mu_\rho=\rho\delta_d+(1-\rho)\delta_a,\qquad B=\rho^{-4}.
\]
Then $R_{\mu_\rho,\rho}(B)\ge\rho/100$, whereas
\[
 H_{\mu_\rho,\rho}(B^{-1/2})
 =\rho^2+4(1-\rho)\rho^{5/2}=O(\rho^2).
\]
The minimax error is therefore of order $\rho$, and the heuristic can underestimate it by an unbounded factor.
\end{theorem}
\begin{proof}
Put $\eta=1/64$ and choose $q_\pm(d)=1/2\pm\eta$, $q_\pm(a)=0$. Both targets lie in $[0,\rho]$ and are separated by $\Delta=2\eta\rho=\rho/32$. Even if the algorithm is told every magnitude, only a queried coin of magnitude $d$ is informative. The original posterior inequality gives, for every next observation and every history,
\[
 \KL_{\rm next}\le\frac{32d^2\eta^2}{(1-4\eta^2)^2}.
\]
The stopped chain rule and $Bd^2=1$ imply $\KL(\mathsf P_+,\mathsf P_-)\le k_0:=32\eta^2/(1-4\eta^2)^2<0.008$. Hence $\TV\le\sqrt{k_0/2}$. Thresholding an estimator halfway between the two targets gives error sum at most $8R^2/\Delta^2$, whereas every test has error sum at least $1-\TV$. Thus
\[
 R\ge\Delta\sqrt{(1-\sqrt{k_0/2})/8}>\rho/100.
\]
The constant estimator $\rho/2$ has error at most $\rho/2$. Finally, $(\rho B)^{-1/2}=\rho^{3/2}$ lies between $d$ and $a$. Direct substitution in \eqref{loc:eq:heuristic} gives the stated expression.
\end{proof}

The hard subpopulation contains both positive and negative near-threshold coins. Its local positive fraction is close to $1/2$, although its total population mass is only $\rho$. Multiplying the global modulus by $\rho$ incorrectly treats the difficult portion as if positives were thinly spread across all magnitudes. Section~\ref{loc:sec:power} also gives a failure example with one fixed, uniform magnitude law, so this is not merely an artifact of changing $\mu$ with $\rho$.

\subsection{Two localized lower bounds for arbitrary laws}
\subsubsection{A balanced-region bound independent of rare-prior machinery}
Define
\begin{equation}\label{loc:eq:Gamma}
 \Gamma_{\mu,\rho}(h)=
 \sup_{A:\,\mu(A)\le\rho}\int_A\min\{1,h/d\}\,\mu(dd).
\end{equation}
\begin{proposition}\label{loc:prop:Gamma}
There is a universal $c>0$ such that
\[
 R_{\mu,\rho}(B)\ge c\Gamma_{\mu,\rho}(B^{-1/2}),\qquad
 r_{\mu,\rho}(B,\delta)\ge c\Gamma_{\mu,\rho}\left(\sqrt{\log(1/\delta)/B}\right).
\]
This lower bound remains valid with revealed or selected magnitudes.
\end{proposition}
\begin{proof}
For any admissible $A$, put
$q_\pm(d)=\ind_A(d)[1/2\pm\eta_h(d)]$ with
$\eta_h(d)=\frac18\min(1,h/d)$. The larger target is at most $5\mu(A)/8\le\rho$. The target separation is
$\Delta_h=\frac14\int_A\min(1,h/d)\,d\mu$.
Every observation has predictive KL at most $128h^2/225$ by \cref{sec:lower}, and observations outside $A$ have zero divergence. Thus transcript KL is at most $128Bh^2/225$.
For RMS choose $h=c_0B^{-1/2}$ with a sufficiently small universal $c_0$ and use the testing argument above. For confidence use $h=c_1\sqrt{\log(1/\delta)/B}$, choosing $c_1$ so that transcript KL is strictly less than $\kl(1-\delta,\delta)$; an estimator with radius smaller than $\Delta_h/2$ would contradict data processing. The integrand satisfies $\min(1,\lambda h/d)\ge\lambda\min(1,h/d)$ for $0<\lambda\le1$. Taking the supremum over $A$ proves the claims without changing the scale.
\end{proof}

Atoms can prevent $\Gamma$ from using a small amount of mass at a magnitude. The next functional permits fractional allocation and includes the fixed-gap rarity factor.

\subsubsection{A variational lower modulus}
Define
\begin{equation}\label{loc:eq:Psi}
 \Psi_{\mu,\rho}(h)=
 \sup_{\substack{0\le w\le1\\\int w\,d\mu\le\rho}}
 \int\min\{w(d),\,h\sqrt{w(d)}/d\}\,\mu(dd).
\end{equation}
The optimization variable $w$ represents how a hard alternative allocates its available positive mass. It is not an assumed or estimated sign function of the true population.

\begin{lemma}[Published single-coin input]\label{loc:lem:LV}
There are universal constants $c_*,C_*>0$ such that, for $0<u\le c_*$, $0\le v\le c_*u$, $0<d\le1/2$, and any stopped single-coin adaptive experiment, its transcript laws under sign priors $u$ and $u+v$ satisfy
\begin{equation}\label{loc:eq:single}
 H^2(\mathsf P_u,\mathsf P_{u+v})\le
 C_*\frac{v^2d^2}{u}\E_u\tau.
\end{equation}
Here $H^2(P,Q)=1-\int\sqrt{dP\,dQ}$ and external, parameter-independent randomization is allowed.
\end{lemma}
\begin{proof}[Source and extension]
For sufficiently small $d$, this is the first expected-work bound in Lee--Valiant Corollary 19~\cite{leevaliant}, with their $\rho,\epsilon,\Delta$ replaced by $u,v,d$. Their result averages over histories and is not a history-uniform predictive-KL claim. Its proof is not reproduced here; it is an explicit external theorem used by this note.
For $d\ge d_*>0$, reveal the sign at the first requested observation. Before that request the experiment has no sign information. By data processing, the Hellinger distance is at most the probability of making the first request times the distance between $\operatorname{Ber}(u)$ and $\operatorname{Ber}(u+v)$. The latter is at most $C v^2/u$ when $u$ and $v/u$ are small. Since $\E_u\tau$ is at least that request probability and $d^2\ge d_*^2$, this proves \eqref{loc:eq:single} on the remaining range, including deterministic coins at $d=1/2$. Increasing a universal constant joins the two ranges.
\end{proof}

\begin{lemma}[Bounded-ratio decomposition]\label{loc:lem:hybrid}
Consider two independent-coin population laws $\nu_0,\nu_1$ whose latent Radon--Nikodym ratio lies in $[1-\kappa,1+\kappa]$ for some fixed $\kappa<1$. For a fully adaptive experiment stopped after at most $n$ observations, let $P_0,P_1$ be its transcript laws. Let $P_i$ be the hybrid law using $\nu_1$ only for coin $i$ and $\nu_0$ for every other coin. Then
\[
 \KL(P_0,P_1)\le C_\kappa\sum_i H^2(P_0,P_i).
\]
The same conclusion holds by finite-prefix limits for almost-surely terminating experiments whenever the right-hand side is uniformly bounded along the prefixes.
\end{lemma}
\begin{proof}
Index coins in order of their first observation, so at most $n$ indices occur. Include the algorithm's independent random seed in the transcript. For a fixed transcript, let $L_i=dP_i/dP_0$. This is the local coin likelihood ratio, marginalized over that coin's latent state, so $1-\kappa\le L_i\le1+\kappa$. Independence of latent coins and the parameter-independent action rule give
\[
 \frac{dP_1}{dP_0}=\prod_iL_i.
\]
This identity holds even though the local sequence lengths were chosen using other coins' outcomes. Also $\E_0L_i=1$. On this compact ratio interval,
$-\log x+(x-1)\le C_\kappa(x-1)^2$, and
$(x-1)^2\le C_\kappa(\sqrt x-1)^2$. Taking expectations and using
$\E_0(\sqrt{L_i}-1)^2=2H^2(P_0,P_i)$ proves the finite-prefix claim. Increasing stopped-transcript sigma fields and the standard monotone limit of relative entropy prove the last assertion. Applying the argument at finite prefixes avoids any unjustified interchange of an unbounded stopping sum and a likelihood derivative.
\end{proof}

This lemma is the bounded-relative-change mechanism of Lee--Valiant Lemma 18, stated for the present magnitude-dependent perturbations. It is not a new general decomposition principle. Unlike an attempted uniform posterior bound, it charges informative but rare histories by their expected frequency.

\begin{theorem}[Arbitrary-law rare-event lower bound]\label{loc:thm:PsiLower}
For the ordinary iid fresh-coin experiment, there is a universal $c>0$ such that
\begin{align}
 R_{\mu,\rho}(B)&\ge c\Psi_{\mu,\rho}(B^{-1/2}),\label{loc:eq:psilowr}\\
 r_{\mu,\rho}(B,\delta)&\ge c\Psi_{\mu,\rho}\left(\sqrt{\log(1/\delta)/B}\right).\label{loc:eq:psilowh}
\end{align}
The proof uses Lemma~\ref{loc:lem:LV}; it does not assert this strengthened bound for an arbitrary selected-magnitude oracle.
\end{theorem}
\begin{proof}
Fix any feasible $w$ in \eqref{loc:eq:Psi} and write $f_h(d)=\min\{w(d),h\sqrt{w(d)}/d\}$. Choose small universal constants $b>0$ and $c_0>0$ with $c_0\ll b$ and $b+c_0\le1$. Define
\[
 q_0(d)=bw(d),\qquad q_1(d)=bw(d)+c_0f_h(d).
\]
Both belong to \eqref{loc:eq:class}. Their target separation is $c_0\int f_h\,d\mu$.
Where $w>0$, put $u=bw$, $v=c_0f_h$. The relative change $v/u$ is uniformly small and
\begin{equation}\label{loc:eq:localcost}
 \frac{v^2d^2}{u}\le \frac{c_0^2}{b}h^2.
\end{equation}
Where $w=0$ the two laws agree. The sign-prior likelihood ratios, and hence the latent coin ratios, are uniformly in $[1-\kappa,1+\kappa]$ for a fixed $\kappa<1$.

At a finite observation prefix, apply Lemma~\ref{loc:lem:hybrid}. To bound its $i$th term, simulate all other coins under $q_0$ as external randomness, and reveal $D_i$ only for the purpose of bounding information. Conditional on $D_i=d$, this is a single-coin stopped experiment with priors $u(d),u(d)+v(d)$. Data processing, Lemma~\ref{loc:lem:LV}, and \eqref{loc:eq:localcost} give
\[
 H^2(P_0,P_i)\le C h^2\E_0\tau_i.
\]
Here the conditional distributions of all simulated coins and the algorithmic seed are the same in both worlds. Summing over coins counts each observation once. Thus
\[
 \KL(P_0,P_1)\le C h^2\E_0(T\wedge n)\le Ch^2B.
\]
Let $n$ tend to infinity. For RMS choose a small constant multiple of $B^{-1/2}$; for confidence choose a small constant multiple of $\sqrt{\log(1/\delta)/B}$. The same two-point testing arguments as Proposition~\ref{loc:prop:Gamma} apply. Finally, $\Psi(\lambda h)\ge\lambda\Psi(h)$ for $0<\lambda\le1$, and taking the supremum over $w$ yields the stated scales.
\end{proof}

\subsubsection{An exact solution of the variational problem}
\begin{proposition}\label{loc:prop:optimizer}
For $0<\rho<1$, let $t=t_{\mu,\rho}>0$ be the unique solution of
\begin{equation}\label{loc:eq:tcut}
 \E_\mu\min\{1,t^2/D^2\}=\rho.
\end{equation}
Then $w_*(d)=\min\{1,t^2/d^2\}$ is an optimizer in \eqref{loc:eq:Psi} for every $h>0$. If $h\ge t$, then $\Psi_{\mu,\rho}(h)=\rho$. If $0<h<t$, then
\begin{equation}\label{loc:eq:psiform}
 \Psi_{\mu,\rho}(h)=F_\mu(h)
 +h\E_\mu[D^{-1};\,h<D\le t]
 +ht\E_\mu[D^{-2};\,D>t].
\end{equation}
These identities allow atoms and singular laws.
\end{proposition}
\begin{proof}
The left side of \eqref{loc:eq:tcut} is continuous, tends to zero at $t=0$, and is strictly increasing until it reaches one. This proves existence and uniqueness at $\rho<1$.
For $h<t$, put $\lambda=h/(2t)<1/2$. For fixed $d$, the concave function
$f_d(w)=\min\{w,(h/d)\sqrt w\}$ is linear until $w=(h/d)^2$, then square-root, on $[0,1]$. Maximizing $f_d(w)-\lambda w$ on $[0,1]$ gives exactly $w_*(d)$: in the square-root portion its stationary point is $t^2/d^2$, and otherwise the maximum is at $1$. Integration and $\int w_*\,d\mu=\rho$ prove optimality. Splitting at $h$ and $t$ gives \eqref{loc:eq:psiform}, with the endpoint contributions agreeing if there are atoms. If $h\ge t$, $f_d(w_*)=w_*$ everywhere, so the objective equals the maximum possible value $\rho$.
\end{proof}

In particular,
\begin{equation}\label{loc:eq:anchorspsi}
 \Psi_{\delta_\Delta,\rho}(h)=\min\{\rho,h\sqrt\rho/\Delta\},\qquad
 \rho\omega_\mu(h/\sqrt\rho)\le\Psi_{\mu,\rho}(h)\le\omega_\mu(h).
\end{equation}
The heuristic is the feasible \emph{constant allocation} $w\equiv\rho$ in the variational problem. Its failure is an allocation failure, not a failure of its two sanity anchors. If $\rho$ stays bounded away from zero, \eqref{loc:eq:anchorspsi} recovers the original modulus up to constants. The inverse-square tail in \eqref{loc:eq:psiform} describes how hard negative coins outside the unresolved core still contribute.

\paragraph{Correction to the earlier conjecture.}
The proposed arbitrary-law matching-$\Psi$ upper bound is refuted by
Theorem~\ref{thm:v5-critical}. The lower bound and optimizer above remain
valid. The matching replacement for arbitrary laws is the
entropy--tail law of \cref{sec:entropy-tail}.

\subsection{Matching rare-event rates for power margins}\label{loc:sec:power}
Let $\mu_{\alpha,a}$ have $F(t)=(t/a)^\alpha$ on $(0,a]$, with $0<a\le1/4$. Write
\begin{equation}\label{loc:eq:rscale}
 \mathcal R_{\alpha,\rho}(x)=
 \begin{cases}
 \min\{\rho,x^\alpha\},&0<\alpha<1,\\
 \min\{\rho,x[1+\logp(\rho/x)]\},&\alpha=1,\\
 \min\{\rho,\rho^{1-1/\alpha}x\},&1<\alpha<2.
 \end{cases}
\end{equation}
\begin{theorem}\label{loc:thm:power}
For fixed $0<\alpha<2$, all $0<a\le1/4$, $0<\rho\le1/4$, $B\ge1$, and $0<\delta\le1/4$,
\begin{align*}
 R_{\mu_{\alpha,a},\rho}(B)&\asymp_\alpha
 \mathcal R_{\alpha,\rho}\bigl((a\sqrt B)^{-1}\bigr),\\
 r_{\mu_{\alpha,a},\rho}(B,\delta)&\asymp_\alpha
 \mathcal R_{\alpha,\rho}\bigl(a^{-1}\sqrt{\log(1/\delta)/B}\bigr).
\end{align*}
The same upper bounds hold for every magnitude law with $F_\mu(t)\le\min\{(t/a)^\alpha,1\}$. The upper construction knows $(\alpha,a,\rho)$, not the individual magnitudes or sign function. It does not give a distribution-blind rare-event algorithm for arbitrary laws.
\end{theorem}

\subsubsection{Lower proof}
Take $A=(0,a\rho^{1/\alpha}]$ in Proposition~\ref{loc:prop:Gamma}. Its mass is $\rho$. Put $t_\rho=a\rho^{1/\alpha}$. For $h\ge t_\rho$ the restricted modulus is $\rho$. For $h<t_\rho$ it is
\begin{equation}\label{loc:eq:restrictedpower}
 \left(\frac h a\right)^\alpha+
 \frac{\alpha h}{a^\alpha}\int_h^{t_\rho}d^{\alpha-2}\,dd.
\end{equation}
For $\alpha\ne1$, the integral is
$(t_\rho^{\alpha-1}-h^{\alpha-1})/(\alpha-1)$; for $\alpha=1$ the result is
$(h/a)[1+\log(a\rho/h)]$. These expressions are comparable, with constants depending only on $\alpha$, to \eqref{loc:eq:rscale}. This proof uses the balanced rare-region construction and does not need the imported single-coin rare-prior lemma.

\subsubsection{Why the original unmodified variance is not rarity-scaled}
\begin{proposition}\label{loc:prop:noise}
Fix an interior magnitude $0<d<1/2$ and let every coin be negative. For the original untruncated geometric single-term estimator $Y$, which is unbiased and square-integrable at this fixed gap,
\[
 \Var Y\ge\frac{1/2-d}{1-2^{-1/2}}>0.
\]
For the truncation at level $L$, the same second-moment lower bound holds, and
$\Var Y_L\ge(1/2-d)/(1-2^{-1/2})-(\E Z_{2^L})^2$.
\end{proposition}
\begin{proof}
The event $J=0$ has probability $p_0=1-2^{-1/2}$ and contributes $Y=X/p_0$ for one Bernoulli observation with success probability $1/2-d$. Therefore $\E Y^2\ge(1/2-d)/p_0$. In the untruncated case $\E Y=0$. In the truncated case telescoping gives $\E Y_L=\E Z_{2^L}$; subtract its square.
\end{proof}
This is an analytic obstruction to deriving a $\rho$ factor by reusing the original second-moment estimate unchanged. At $d=1/4$ the lower floor is approximately $0.8536$, even though $\theta=0$.

\subsubsection{A shifted-start construction}
Retain the original classifier $Z_n$ with dyadic checkpoints and threshold
$\sqrt{[1+\log(n/k)]/k}$ at checkpoint $k<n$. Define
\[
 e_n(d)=\min\{1,(nd^2)^{-2}\},\qquad
 m_n(d)=\min\{n,d^{-2}[1+\logp(nd^2)]\}.
\]
The original conditional inequalities are
$\PP(Z_n\ne S\mid d,S)\le e_n(d)$ and
$\E(\tau_n\mid d,S)\le17m_n(d)$ (\cref{sec:classifier,sec:minimax-proof}).

\begin{lemma}\label{loc:lem:integrals}
For $na^2\ge1$ and $0<\alpha<2$,
\begin{align}
 e_n:=\E_{\mu_{\alpha,a}}e_n(D)&\le K_e a^{-\alpha}n^{-\alpha/2},&K_e&=\frac4{4-\alpha},\label{loc:eq:eb}\\
 \E_{\mu_{\alpha,a}}m_n(D)&\le K_m a^{-\alpha}n^{1-\alpha/2},&
 K_m&=1+\frac\alpha{2-\alpha}+\frac{2\alpha}{(2-\alpha)^2}.\label{loc:eq:mb}
\end{align}
The same bounds hold under the stated stochastic-dominance envelope.
\end{lemma}
\begin{proof}
Split each integral at $d=n^{-1/2}$. For \eqref{loc:eq:eb}, the first part is $(a\sqrt n)^{-\alpha}$; extending the remaining power integral to infinity gives the factor $1+\alpha/(4-\alpha)$. For \eqref{loc:eq:mb}, substitute $u=d\sqrt n$ in the tail and extend its upper endpoint to infinity. The remaining factor is
\[
 1+\alpha\int_1^\infty u^{\alpha-3}(1+2\log u)\,du
 =K_m.
\]
Both integrands are nonincreasing in $d$, which proves the envelope statement. The restriction $\alpha<2$ is essential to this particular work bound.
\end{proof}

Choose the smallest dyadic $n_0$ satisfying
\begin{equation}\label{loc:eq:n0}
 n_0\ge a^{-2}(K_e/\rho)^{2/\alpha}.
\end{equation}
If $B<4n_0$, return $\rho/2$. Otherwise choose $N=n_02^L$ as the largest such number with $4N\le B$. Thus $B/8<N\le B/4$. Set $n_j=n_02^j$ and
\[
 \Delta_0=Z_{n_0},\qquad \Delta_j=Z_{n_j}-Z_{n_{j-1}},\quad j\ge1,
\]
using a shared observation prefix within each increment. Sample $J$ with
$p_j=(1-r)r^j$, $r=2^{-1/2}$, and return $Y_L=\Delta_J/p_J$ if $J\le L$ and zero otherwise. Coins are independent between calls. The observed work $C_L$ is at most $N$.

\begin{lemma}\label{loc:lem:shiftbounds}
With $S_L=\sum_{j=0}^L n_j^{(1-\alpha)/2}$, uniformly over the rare-event envelope class,
\begin{align}
 |\E Y_L-\theta|&\le K_e a^{-\alpha}N^{-\alpha/2},\label{loc:eq:shiftbias}\\
 \E Y_L^2&\le C_\alpha a^{-\alpha}n_0^{-1/2}S_L,\label{loc:eq:shiftvar}\\
 \E C_L&\le C_\alpha a^{-\alpha}n_0^{1/2}S_L.\label{loc:eq:shiftwork}
\end{align}
\end{lemma}
\begin{proof}
Telescoping proves \eqref{loc:eq:shiftbias}. At the base,
$\E Z_{n_0}^2=\PP(Z_{n_0}=1)\le\theta+e_{n_0}\le2\rho$.
Dyadic rounding in \eqref{loc:eq:n0} gives
$\rho\le2^{\alpha/2}K_e a^{-\alpha}n_0^{-\alpha/2}$, so the base obeys the same power envelope as the increments. For $j\ge1$,
\[
 \E\Delta_j^2\le e_{n_j}+e_{n_{j-1}}\le C_\alpha a^{-\alpha}n_j^{-\alpha/2}.
\]
Divide by $p_j=(1-r)\sqrt{n_0/n_j}$ and sum to obtain \eqref{loc:eq:shiftvar}. For work, a coupled increment costs at most the sum of the two classifier costs, so its expectation is at most $34K_m a^{-\alpha}n_j^{1-\alpha/2}$; the base needs only the factor $17$. Multiply by $p_j$ and sum to obtain \eqref{loc:eq:shiftwork}.
\end{proof}

\subsubsection{Calibration and completion of the upper proof}
Apply the original two-success wrapper to iid pairs $(Y_L,C_L)$, with cap $N$. Its pilot markers have success probability $C_L/N$. Stop at the second success, at pilot count $K$, discard the pilot values, and average $K-1$ fresh production values. With $c=\E C_L>0$ and $v=\Var Y_L$, the identities are
\begin{equation}\label{loc:eq:calibration}
 \E T=4N-c\le B,\quad \E\widehat\theta=\E Y_L,\quad
 \Var\widehat\theta=vc/N.
\end{equation}
The identity $\E[1/(K-1)]=c/N$ is classical inverse-binomial estimation; the wrapper is inherited from \cref{sec:minimax-proof}, not claimed as new~\cite{haldane,mendo}. Independence of the fresh production samples is essential.

Equations \eqref{loc:eq:shiftbias}--\eqref{loc:eq:shiftwork} imply
\begin{equation}\label{loc:eq:upperwork}
 \bigl(\E|\widehat\theta-\theta|^2\bigr)^{1/2}
 \le C_\alpha\left[a^{-\alpha}N^{-\alpha/2}
       +a^{-\alpha}N^{-1/2}S_L\right].
\end{equation}
For $\alpha<1$, $S_L\le C_\alpha N^{(1-\alpha)/2}$. For $\alpha=1$, $S_L=L+1$. For $\alpha>1$, $S_L\le C_\alpha n_0^{(1-\alpha)/2}$. Substitution of \eqref{loc:eq:n0} and $N\asymp B$ yields \eqref{loc:eq:rscale}. In the small-budget branch, the right side of the target rate is at least a constant multiple of $\rho$, so returning $\rho/2$ suffices. Projecting the final estimate onto $[0,\rho]$ cannot increase squared error.

For high confidence, use an odd number $k\asymp\log(1/\delta)$ of independent groups, each with expected budget $B/k$, and take their median. Chebyshev gives failure at most $1/8$ at $\sqrt8$ times each group's RMS bound; a binomial tail bound gives median failure at most $e^{-9k/32}$. The rate follows after replacing $B$ by $B/k$. If group budgets are too small, the constant estimator handles the saturated regime. Expected costs add. This completes Theorem~\ref{loc:thm:power}.

\subsubsection{Consequences and limits of the theorem}
For $0<\varepsilon\le c_\alpha\rho$, the required expected work is, up to $\alpha$-dependent factors,
\begin{center}\begin{tabular}{ll}\toprule
Margin & Work, multiplied by $a^{-2}\log(1/\delta)$\\\midrule
$0<\alpha<1$ & $\varepsilon^{-2/\alpha}$\\
$\alpha=1$ & $\varepsilon^{-2}[1+\log(\rho/\varepsilon)]^2$\\
$1<\alpha<2$ & $\rho^{2-2/\alpha}\varepsilon^{-2}$\\\bottomrule
\end{tabular}\end{center}
At constant $\rho$ these reduce to the corresponding original rates. The fixed-gap rate \eqref{loc:eq:lvanchor} is a separate cited special case; it is not obtained by extrapolating Theorem~\ref{loc:thm:power} beyond its hypotheses.

For the fixed uniform law ($\alpha=1$), set $x=(a\sqrt B)^{-1}=\rho^2$. The proved rate is of order $\rho^2\log(1/\rho)$, whereas the proposed heuristic is of order $\rho^{5/2}\log(1/\rho)$. Thus failure of the heuristic also occurs for a single fixed $\mu$.

The optimizer formula predicts an additional boundary at $\alpha=2$: in a sufficiently large-budget regime its lower modulus has scale $(h/a)\sqrt{\rho\log(1/\rho)}$, versus $(h/a)\sqrt\rho$ for $\alpha>2$. The above-two rate is matched in \cref{app:above-two}; at the quadratic
law, \cref{app:critical} proves the additional scale-entropy factor. Constants in the proved theorem need not be uniform as $\alpha$ approaches $1$ or $2$.

\subsection{A certified general upper bracket, not a closed characterization}
For any finite dyadic sequence $n_0,\ldots,n_L=N$ with $4N\le B$, let
\[
 e_n=\E_\mu e_n(D),\quad
 v_0=\rho+e_{n_0},\quad v_j=e_{n_j}+e_{n_{j-1}},\quad
 b_0=17\E_\mu m_{n_0},\quad b_j=34\E_\mu m_{n_j}\ (j\ge1).
\]
All these quantities are finite at finite caps. Put
$S=\sum_{j=0}^L\sqrt{v_jb_j}$ and choose a finite categorical law
$p_j\propto\sqrt{v_j/b_j}$. Sample $\Delta_J/p_J$ and apply \eqref{loc:eq:calibration} with cap $N$.
\begin{proposition}\label{loc:prop:generalupper}
This estimator has expected work at most $B$ and RMS error at most
\[
 \sqrt{e_N^2+S^2/N}.
\]
Consequently the minimax risk is bounded above by the infimum of this expression over the finite grids, and also by $\rho/2$ and by the clipped original upper bound. The allocation here requires the stated integrals of $\mu$ or valid upper bounds for them.
\end{proposition}
\begin{proof}
The same coupling argument gives $\E\Delta_j^2\le v_j$ and $\E W_j\le b_j$. Thus
$\E Y^2\le\sum v_j/p_j$ and $\E C\le\sum p_jb_j$. Their product equals at most $S^2$ for the stated choice of probabilities. Telescoping bounds the bias by $e_N$, and calibration bounds the variance by $S^2/N$.
\end{proof}

This bracket is constructive, but it does not remove the fixed-gap screening difficulty or establish arbitrary-law optimality. The sharp arbitrary-law replacement is \cref{thm:entropy-tail-law};
the matching-$\Psi$ upper objective is false by \cref{thm:v5-critical}. It would be incorrect to present the variational identity as if it supplied that algorithm.
